\chapter*{Abstract}

Norwegian transport companies operate without systematic visibility into resource utilisation: overtime is handled reactively, idle time between assignments is invisible in current tools, and load balancing depends on the traffic coordinator's memory. This thesis asks how an algorithm-assisted resource planning platform can improve resource utilisation in Norwegian transport companies while remaining accountable to the coordinator who operates it.

The thesis follows Design Science Research, applied through Peffers' six DSRM activities and six named iterations. The empirical foundation is semi-structured interviews with seven traffic coordinators and transport managers across Norway, analysed through Braun and Clarke's six-phase thematic analysis. The artefact, Ressursplanlegger (RP), is a Next.js + tRPC + PostgreSQL platform with three optimisation engines (a greedy heuristic, OR-Tools CP-SAT, and Timefold) and a Gantt-style drag-and-drop human-in-the-loop timeline. The evaluation framework benchmarks the engines against three synthetic datasets and is complemented by a requirements traceability matrix.

The interviews surface a resource-utilisation visibility gap as the primary finding and an operator-vs-owner asymmetry that mirrors Bainbridge's \emph{Ironies of Automation}: demand for automation is articulated by owners and by Admmit, not by the coordinators who must run any such system. The discussion organises primary findings under three locked anchor concepts: \textbf{Efficiency} (visibility gap, asymmetry, multi-engine benchmark), \textbf{Control} (four-layer human-in-the-loop theory applied to override authority), and \textbf{Adaptability} (configurable soft-constraint weights across companies). Twelve hierarchical limitations (L1 through L12) bound what the thesis can claim, the most consequential being the absence of production deployment (L6) and of user testing with coordinators (L8). The closing claim is that algorithm-assisted planning under stakeholder asymmetry in Norwegian transport is feasible when the operator keeps authority over every algorithm-generated assignment, when the system explains itself in the operator's vocabulary rather than in solver internals, and when what distinguishes one company from another lives as configuration rather than as code.
